\section{Introduction}

Two language models can solve the same exercise correctly while teaching very
differently. One gives the answer and explains it; another asks the student to
complete a step; a third acknowledges the student's frustration before doing
either. Such differences motivate a question about \emph{educational
personality}: do deployed models exhibit recurring teaching tendencies that help
predict how they will respond in a new educational situation?

Answering this question requires more than distinguishing the models' text.
Length, formatting, task instructions, and access to the solution can all produce
recognizable differences. Nor does a model need to behave identically in every
situation to have a recurring tendency. A tutor may consistently explain more
after an erroneous attempt, or acknowledge difficulty specifically when a student
expresses frustration. The empirical question is whether a model's defaults or
its conditional response patterns recur sufficiently to predict new behavior.
An observed difference, a repeatable measurement, and a successful prediction
are distinct steps in that argument.

Education provides concrete choices through which to test this question without
assuming a human personality taxonomy. Revealing a solution and inviting the
student to reason are observable actions; they can co-occur, and neither is
universally the better teaching move. Acknowledging expressed frustration is
also observable without inferring that the model experiences empathy. These
behaviors let us distinguish how a model tends to teach from how accurately it
solves a problem or how much a learner benefits. Our object of study is the
\emph{tutor model's} recurring behavior. It differs from adapting instruction to
an assigned or inferred \emph{student personality}, as in personality-aware
teaching systems \citep{rooein2026pats}.

Prior studies identify model differences through personality-oriented
scenarios and situated behavioral probes \citep{lee2025trait,liu2026bma}.
Educational studies already compare characteristic tutoring-action mixtures
and learn tutor-persona variation from dialogue
\citep{kucheria2025cima,lee2026tutorpersonas}. Building on these observations,
we test the incremental predictive information in deployed models' defaults
and student-state response patterns. The comparison is against the educational
situation, model-specific subject habits, and simple training-output style
profiles, with forecasts locked before new responses are generated.

We organize the investigation around three questions. \textbf{RQ1:} Under the
same educational input, do models show distinguishable defaults that recur
across repeated generations? \textbf{RQ2:} Can those defaults, or model-specific
responses to student state, predict actions on independent source problems?
\textbf{RQ3:} How do the resulting profiles change under neutral rewordings of
the tutor role and explicit teaching instructions? These questions form one
evidence chain. The first establishes recurring differences; the second tests
their predictive value; the third identifies conditions under which they persist
or move.

The study connects an existing educational-output archive to a separate
prospective experiment. The archive contains more than one million raw
prediction records, with different roles, coverage, and repetition histories.
We audit those distinctions before using the archive for exploration. We then
develop observable event measures on a small, permanently excluded pilot and
freeze a source-separated training and confirmation design. Predictions are
fitted and saved before confirmation responses are generated. Model-by-subject
baselines prevent subject-specific habits from being mistaken for student-state
response patterns, while neutral role wording and explicit teaching policies
are manipulated separately on the same inputs.

The resulting claims concern observable tendencies of specified deployments.
They do not establish human-equivalent traits, a latent psychological mechanism,
or learning gains. In particular, a successful default profile does not imply
that a more elaborate conditional profile will improve prediction, and a large
instruction effect does not by itself establish or refute educational
personality. The final conclusions must follow these comparisons, including
measurement failures and limits on transfer.
